\section{Résultats et discussion}
\label{sec:resultats}

Cette section présente les résultats obtenus pour chaque architecture évaluée
lors de la Phase~1 (benchmarking). Pour chaque modèle, nous reportons les
courbes d'entraînement (perte et précision), la matrice de confusion sur
l'ensemble de test, ainsi que le rapport de performance détaillé
(Accuracy, Recall, F1-Score, AUC-ROC, et métriques par classe).

%====================================================================
\subsection{Baseline : CNN personnalisé (\texttt{custom\_cnn})}
\label{subsec:baseline_custom_cnn}

\begin{figure}[H]
    \centering
    \figIfExists[\linewidth]{./plots/training_history_baseline_custom_cnn.png}
    \caption{Courbes d'entraînement --- \texttt{baseline / custom\_cnn}.}
    \label{fig:hist_baseline_custom_cnn}
\end{figure}

\begin{figure}[H]
    \centering
    \figIfExists[0.85\linewidth]{./plots/confusion_matrix_baseline_custom_cnn.png}
    \caption{Matrice de confusion --- \texttt{baseline / custom\_cnn}.}
    \label{fig:cm_baseline_custom_cnn}
\end{figure}

\inputIfExists{./tab/perf_report_baseline_custom_cnn.tex}

\textit{Discussion : à compléter (comportement de la convergence, classes
les plus confondues, écart train/val, etc.).}

%====================================================================
\subsection{Baseline : ResNet-18 (\texttt{resnet18})}
\label{subsec:baseline_resnet18}

\begin{figure}[H]
    \centering
    \figIfExists[\linewidth]{./plots/training_history_baseline_resnet18.png}
    \caption{Courbes d'entraînement --- \texttt{baseline / resnet18}.}
    \label{fig:hist_baseline_resnet18}
\end{figure}

\begin{figure}[H]
    \centering
    \figIfExists[0.85\linewidth]{./plots/confusion_matrix_baseline_resnet18.png}
    \caption{Matrice de confusion --- \texttt{baseline / resnet18}.}
    \label{fig:cm_baseline_resnet18}
\end{figure}

\inputIfExists{./tab/perf_report_baseline_resnet18.tex}

\textit{Discussion : à compléter (apport du préentraînement ImageNet,
gain par rapport au CNN \emph{from scratch}, etc.).}

%====================================================================
\subsection{Transformer : ViT-B/16 (\texttt{vit\_b\_16})}
\label{subsec:transformers_vit_b_16}

\begin{figure}[H]
    \centering
    \figIfExists[\linewidth]{./plots/training_history_transformers_vit_b_16.png}
    \caption{Courbes d'entraînement --- \texttt{transformers / vit\_b\_16}.}
    \label{fig:hist_transformers_vit_b_16}
\end{figure}

\begin{figure}[H]
    \centering
    \figIfExists[0.85\linewidth]{./plots/confusion_matrix_transformers_vit_b_16.png}
    \caption{Matrice de confusion --- \texttt{transformers / vit\_b\_16}.}
    \label{fig:cm_transformers_vit_b_16}
\end{figure}

\inputIfExists{./tab/perf_report_transformers_vit_b_16.tex}

\textit{Discussion : à compléter (efficacité de l'attention globale,
sensibilité à la taille du jeu de données, etc.).}

%====================================================================
\subsection{Transformer : Swin-T (\texttt{swin\_t})}
\label{subsec:transformers_swin_t}

\begin{figure}[H]
    \centering
    \figIfExists[\linewidth]{./plots/training_history_transformers_swin_t.png}
    \caption{Courbes d'entraînement --- \texttt{transformers / swin\_t}.}
    \label{fig:hist_transformers_swin_t}
\end{figure}

\begin{figure}[H]
    \centering
    \figIfExists[0.85\linewidth]{./plots/confusion_matrix_transformers_swin_t.png}
    \caption{Matrice de confusion --- \texttt{transformers / swin\_t}.}
    \label{fig:cm_transformers_swin_t}
\end{figure}

\inputIfExists{./tab/perf_report_transformers_swin_t.tex}

\textit{Discussion : à compléter (apport des fenêtres décalées,
hiérarchie multi-échelle, etc.).}

%====================================================================
\subsection{Hybride : MaxViT-T (\texttt{maxvit\_t})}
\label{subsec:hybrids_maxvit_t}

\begin{figure}[H]
    \centering
    \figIfExists[\linewidth]{./plots/training_history_hybrids_maxvit_t.png}
    \caption{Courbes d'entraînement --- \texttt{hybrids / maxvit\_t}.}
    \label{fig:hist_hybrids_maxvit_t}
\end{figure}

\begin{figure}[H]
    \centering
    \figIfExists[0.85\linewidth]{./plots/confusion_matrix_hybrids_maxvit_t.png}
    \caption{Matrice de confusion --- \texttt{hybrids / maxvit\_t}.}
    \label{fig:cm_hybrids_maxvit_t}
\end{figure}

\inputIfExists{./tab/perf_report_hybrids_maxvit_t.tex}

\textit{Discussion : à compléter (complémentarité conv + attention,
coût computationnel comparé à un ViT pur, etc.).}

%====================================================================
\subsection{SSM : Vision Mamba (\texttt{vision\_mamba})}
\label{subsec:ssm_vision_mamba}

\begin{figure}[H]
    \centering
    \figIfExists[\linewidth]{./plots/training_history_ssm_vision_mamba.png}
    \caption{Courbes d'entraînement --- \texttt{ssm / vision\_mamba}.}
    \label{fig:hist_ssm_vision_mamba}
\end{figure}

\begin{figure}[H]
    \centering
    \figIfExists[0.85\linewidth]{./plots/confusion_matrix_ssm_vision_mamba.png}
    \caption{Matrice de confusion --- \texttt{ssm / vision\_mamba}.}
    \label{fig:cm_ssm_vision_mamba}
\end{figure}

\inputIfExists{./tab/perf_report_ssm_vision_mamba.tex}

\textit{Discussion : à compléter (complexité linéaire vs. quadratique
de l'attention, stabilité de l'entraînement, etc.).}

%====================================================================
\subsection{KAN : Vision KAN (\texttt{vision\_kan})}
\label{subsec:kan_vision_kan}

\begin{figure}[H]
    \centering
    \figIfExists[\linewidth]{./plots/training_history_kan_vision_kan.png}
    \caption{Courbes d'entraînement --- \texttt{kan / vision\_kan}.}
    \label{fig:hist_kan_vision_kan}
\end{figure}

\begin{figure}[H]
    \centering
    \figIfExists[0.85\linewidth]{./plots/confusion_matrix_kan_vision_kan.png}
    \caption{Matrice de confusion --- \texttt{kan / vision\_kan}.}
    \label{fig:cm_kan_vision_kan}
\end{figure}

\inputIfExists{./tab/perf_report_kan_vision_kan.tex}

\textit{Discussion : à compléter (apport des fonctions de Kolmogorov-Arnold,
expressivité par paramètre, etc.).}

%====================================================================
\subsection{Synthèse comparative}
\label{subsec:synthese}

\textit{Discussion globale à compléter : tableau récapitulatif
(Accuracy, F1, AUC-ROC, \#params, FLOPs) sur l'ensemble des modèles,
choix de l'architecture retenue pour la Phase~2 et justification
(compromis précision / coût / généralisation).}